\section{Commutative Lemma}

In this section, we discuss the formalization of some commutativity results from the previous
chapter. Fix a formal group law $F(X,Y)$ over $R$ throughout this section.
First, we fix some ambient variables.

\begin{lstlisting}
variable {R : Type*} [CommRing R] (F : FormalGroup R) [Nontrivial R] {σ : Type*}
\end{lstlisting}

According to the definition of commutativity for formal group laws, we have that $F(X,Y)$ is
commutative if and only if the coefficients of $X^i Y^j$ and $X^j Y^i$ coincide.

\begin{lstlisting}
theorem comm_iff_coeff_symm :
  F.comm ↔ ∀ (i j : ℕ), coeff (single 0 i + single 1 j) F.toFun
    = coeff (single 0 j + single 1 i) F.toFun := by
\end{lstlisting}

A precommutator of a formal group law \lstinline|F| is a multivariate power series
\lstinline|X +[F] Y +[F] (-X)|. Typically, a commutator of a group law is defined to be
$x y x^{-1} y^{-1}$, where $x, y$ are elements of a group $G$, and a group is commutative if and only
if its commutator is zero.
A formal group law $F(X,Y)$ is commutative if and only if the precommutator is equal to
\lstinline|Y|.

\begin{lstlisting}
def preCommutator (f g : MvPowerSeries σ R) := f +[F] g +[F] (addInv F f)

def commutator : MvPowerSeries (Fin 2) R :=
    X₀ +[F] X₁ +[F] (addInv F X₀) +[F] (addInv F X₁)

lemma preCommutator_ne_of_nonComm (h : ¬ F.comm) :
    preCommutator F X₀ X₁ ≠ X₁ := by

lemma comm_iff_commutator_eq_zero : F.comm ↔ commutator F = 0 := by
\end{lstlisting}

Now we can formalize Lemma \ref{lem: non_com_hom} and Lemma \ref{lem: target_com}.

\begin{lstlisting}
theorem exists_nonzero_hom_to_Ga_or_Gm_of_not_comm (h : ¬ F.comm) :
    (∃ (α : FormalGroupHom F (Gₐ (R := R))), α.toFun ≠ 0) ∨
    (∃ (α : FormalGroupHom F (Gₘ (R := R))), α.toFun ≠ 0) :=

theorem comm_of_exists_nonzero_hom_to_comm (F' : FormalGroup R) [IsDomain R]
    (α : FormalGroupHom F F') (ha : α.toFun ≠ 0) :
    F'.comm → F.comm := by
\end{lstlisting}

The instance \lstinline|IsDomain R| means that \lstinline|R| is an integral domain.
Here \lstinline|Gₐ| and \lstinline|Gₘ| are abbreviations for the additive and multiplicative formal group laws.
The notation \lstinline|(R := R)| means that the coefficient ring of the
formal group laws is \lstinline|R|.

Therefore, we can deduce that any formal group law over a domain $R$ is commutative, with a
neat proof.

\begin{lstlisting}
theorem comm_of_isDomain [IsDomain R] : F.comm := by
  by_contra hc
  obtain ⟨α, ha⟩| ⟨α, ha⟩ := exists_nonzero_hom_to_Ga_or_Gm_of_not_comm F hc
  · exact hc ((comm_of_exists_nonzero_hom_to_comm _ _ α ha) Gₐ.comm)
  · exact hc ((comm_of_exists_nonzero_hom_to_comm _ _ α ha) Gₘ.comm)
\end{lstlisting}

